\section{Limitations and Threats to Validity}\label{sec:limits}

\paragraph{Parallel trends.} The most significant threat to identification is the violation of the parallel trends assumption documented in Section~\ref{sec:results-pretrend}. The full-panel event study shows a V-shaped pre-treatment divergence between treated and never-treated tracts: group-time ATTs reach $-\$16$/month at $t = -8$ before recovering to zero at $t = -1$, and a joint Wald test rejects the flat pre-trend null at $p < 0.001$. This divergence reflects the endogenous targeting of prohibitions: high-rent, fast-appreciating neighborhoods were most likely to adopt STR bans, so the treated group had systematically different counterfactual rent trajectories than the control group.

The standard response to a pre-trend violation --- controlling for pre-treatment trends --- introduces its own problems here. Linear tract-specific detrending extrapolates the late recovery phase of the V-shaped pre-period forward and over-subtracts approximately \$15/month by the end of the post-period, mechanically inverting the sign of the residualized ATT (Appendix~\ref{app:residualized}). This is not a real negative effect; it is a functional form artifact. Quadratic detrending partially corrects the curvature misfit but does not yield a rent-comparable point estimate.

Our response is to report HonestDiD sensitivity bounds rather than claim point identification. Following \citet{rbs2023}, we parameterize the violation by a smoothness constraint $M$ capturing the maximum allowable month-to-month change in the counterfactual trend relative to the observed pre-period path. Using the empirical second differences of the pre-period TWFE coefficients, the largest observed value is $M = 2.71$ \$/month per period. Under this bound, the identified set for the average ATT remains consistent with positive effects throughout the post-period window. We interpret this as evidence that the sign of the effect is robust to violations of the magnitude actually present in the data, even if the precise magnitude is not point-identified. The pretrend diagnostic table in Appendix~\ref{app:parallel-trends} reports these heuristics; full Rambachan--Roth sensitivity plots require the \texttt{HonestDiD} \textsc{R} package applied to our group-time ATT output and are documented as a planned extension.

Furthermore, \citet{roth2022pretest} documents that pre-trend tests routinely fail to detect economically relevant violations in settings with limited pre-treatment variation, and that passing a pre-test should not be treated as validation of parallel trends. Our TWFE regression fails to reject at $p \approx 0.34$ despite large and statistically significant pre-period gaps in the CS event study. This divergence between the slope test and the joint event-study test reflects the well-known low power of the former: the TWFE test has essentially no power against the V-shaped alternative observed here.

\paragraph{Matching and sample restriction.} Five-feature $k$-nearest-neighbor matching restricts the analysis to 274 treated tracts and 176 controls (450 of 842 rent-panel tracts). Post-match standardized mean differences improve substantially on slope (from 1.36 to 0.75) but remain material, suggesting that the matched sample does not fully resolve the pre-trend imbalance. The matched estimand is local: it recovers the ATT for the subset of treated tracts that have sufficiently similar never-treated counterparts, not for the full universe of ever-treated tracts. If the treated tracts excluded from matching (due to poor pre-trend matches) have systematically different treatment effects, the matched ATT may not generalize to the full treated population.

Additionally, the three lagged rent features used in matching are computed relative to a panel-wide anchor (the first month of treatment in the full panel) rather than relative to each tract's own treatment date. For late-adopting cohorts, this means the lag features are computed many months before treatment begins, potentially missing trajectory changes in the intervening period. This is a known limitation of the current matching implementation and a priority for future refinement.

\paragraph{Measurement error in the outcome.} The ZORI outcome is a ZIP-level index, disaggregated to census tracts via area-weighted interpolation. This introduces classical measurement error: the rent signal for a tract is an area-weighted average of its overlapping ZIP codes' indices, which smooths over within-ZIP variation and introduces attenuation bias. Tracts with large footprints relative to their ZIP code coverage --- common in the outer suburban portions of Cook County --- have less precise rent signals than tracts in the dense urban core where ZIP and tract boundaries more closely coincide. This measurement error attenuates our estimates toward zero, so our reported ATTs are conservative in this dimension. We do not apply additional deseasonalization beyond Zillow's seasonal adjustment; the seasonality diagnostic (Appendix~\ref{app:seasonality}) confirms that calendar-month fixed effects leave the ATT essentially unchanged (\$5.456 vs.\ \$5.462/month).

\paragraph{SUTVA and spatial spillovers.} The difference-in-differences design assumes that the treatment of one tract does not affect control tracts' outcomes. Section~\ref{sec:results-spillovers} documents evidence against this assumption: never-treated tracts adjacent to treated tracts show substantially lower estimated ATTs ($+\$0.5$/month) than isolated never-treated tracts ($+\$26.0$/month) when used as controls on the five-feature matched panel. This pattern is consistent with positive spatial spillovers from prohibition activity to adjacent neighborhoods, which contaminate the adjacent never-treated controls and attenuate the headline ATT estimate. Our five-feature matched ATT of \$5.46/month (two-feature baseline \$4.60/month) should therefore be interpreted as a lower bound on the direct treatment effect among the studied tracts; the full policy impact, accounting for spillovers to adjacent neighborhoods, is larger.

\paragraph{External validity.} Our results apply to census tracts in Cook County covered by Zillow's rent index that experienced meaningful STR prohibition activity under Chicago's ordinance between 2015 and 2024. The generalizability of these findings to other cities depends critically on the institutional features of the prohibition mechanism. Chicago's building-level opt-in prohibition differs from city-wide platform bans (such as New York City's 2023 restrictions), licensing caps (as in New Orleans and San Francisco), and neighborhood-level zoning restrictions (as in Barcelona). Whether demand-signaling operates similarly across these institutional forms is an open empirical question.

\paragraph{Causal mechanism.} Our design identifies the reduced-form rent effect of crossing the treatment threshold, but cannot separately identify the supply-side and demand-signaling channels discussed in Section~\ref{sec:policy}. The positive sign of our ATT is inconsistent with a pure supply-release channel (which predicts lower rents) and more consistent with a demand-signaling channel, but we cannot rule out confounding from neighborhood upgrading or gentrification acceleration that coincides with prohibition adoption. Individual-level tenant data linking residential moves to building prohibition status would be needed to distinguish these mechanisms.
